\subsection{A stronger bound from a fixed-norm element}\label{sec:fixed-norm}
For the height and kernel estimates in this subsection use only
\eqref{eq:packet1}, and impose the additional conditions that $A,B$ are
squarefree and
\[
 A\equiv22\pmod {24},\qquad B\equiv23\pmod {24}.
\]
They specify the squarefree branch under consideration; the more general
packet estimates above did not require them. No Chebyshev decomposition
is needed except in the displayed consequence for $p$. In particular $b\ge22$,
$A,B>1$, $3\nmid AB$, $\gcd(A,B)=1$, and $D=3AB$ is squarefree.
Let $R_A,R_B$ denote the regulators of the maximal real quadratic orders.

We use two general tools from the proof of the new result
\cite[Theorem~2.4 and Lemma~2.5]{PastenNew}, without applying its
irreducible-cubic theorem to a split polynomial. First, in a totally real
degree-four field, for a multiplicative group generated modulo torsion
by three independent elements $g_1,g_2,g_3$ and $g\ne1$ in that group,
\begin{equation}\label{eq:group-approximation}
 -\log|1-g|\le C_*\log\max\{2,\height(g)\}
                  \prod_{i=1}^3\max\{1,\height(g_i)\},
\end{equation}
where $C_*>0$ is effective and absolute. Second, for an algebraic integer
$\omega$ of absolute norm $q>1$ in a real quadratic field $K$, the group
$\OO_K^\times\langle\omega\rangle$ has independent elements $\xi_1,\xi_2$
generating a subgroup modulo torsion of index $1\le N\le2$ such that
\begin{equation}\label{eq:small-generators}
 \omega^N=\zeta\xi_1^{e_1}\xi_2^{e_2},\quad\zeta\in\{\pm1\},
 \qquad \height(\xi_1)\height(\xi_2)\le\frac38R_K\log q.
\end{equation}
Here the set $S$ consists of both real places and no finite place, so the
all-or-none condition in Lemma~2.5 is satisfied and $q'=q$.
The small-generator estimate follows from the exterior-product theorem
of Akhtari--Vaaler \cite[Theorem~1.2]{AV}; its joint use of the unit and
nonunit generators is what saves a regulator factor.
Finally Landau's fixed-degree estimate gives an effective absolute $C_Q$
with
\begin{equation}\label{eq:quadratic-regulator}
 R_A\le C_Q\sqrt A\log(2A),\qquad R_B\le C_Q\sqrt B\log(2B).
\end{equation}

\begin{theorem}\label{thm:fixed-norm-A}
There is an effective absolute constant $C_A$ such that every packet in
this subsection satisfies
\[
 H\le C_A\sqrt A[\log(2A)]^2.
\]
\end{theorem}
\begin{proof}
Put $\varphi=(1+\sqrt5)/2$ and $\rho=\tfrac12\log\varphi$.
Every nonzero algebraic number of degree at most two which is not a root
of unity has height at least $\rho$. Indeed a rational such number has
height at least $\log2$. For a quadratic number, its primitive minimal
polynomial has Mahler measure at least two if the leading or constant
coefficient has absolute value at least two. Otherwise it is a quadratic
unit. For norm minus one a nonzero integral trace gives a root of absolute
value at least $\varphi$; trace zero is reducible. For norm plus one,
traces of absolute value at most two give roots of unity or reducibility,
and the other traces give a root of absolute value at least $\varphi^2$.
The formula $2\height=\log M$ proves the height lower bound.

In $K=\Q(\sqrt A)$ take
\[
 \omega=s+u\sqrt A=\sqrt{b+3}+\sqrt b,
 \qquad \eta=s+r_0\sqrt3=\sqrt{b+3}+\sqrt{b+2}.
\]
Their respective quadratic norms are three and one. The latter is a
positive power of $\varepsilon_3=2+\sqrt3$. Apply
\eqref{eq:small-generators} with $q=3$. The field
$F=\Q(\sqrt A,\sqrt3)$ has degree four, since $A>1$ is squarefree and
$A\ne3$. Define $\gamma=\omega/\eta$. At the chosen real embedding,
\[
 0<\gamma<1,\qquad
 1-\gamma=
 \frac{2}{(\sqrt{b+2}+\sqrt b)(\sqrt{b+3}+\sqrt{b+2})}
 <\frac1{2b}.
\]
As $1\le N\le2$, it follows that $0<1-\gamma^N<1/b$.

The group $\langle-1,\xi_1,\xi_2,\varepsilon_3\rangle$ has rank three
modulo torsion and contains $\gamma^N$. Independence follows because
$K\cap\Q(\sqrt3)=\Q$, and a rational power of $\varepsilon_3$ must equal
its conjugate, forcing that power to be zero. The other two generators
are independent by construction. Moreover $\height(\varepsilon_3)<1$ and
$\max(1,\height(\xi_i))\le\height(\xi_i)/\rho$, giving
\[
 \prod_i\max(1,\height(g_i))\le\frac{3\log3}{8\rho^2}R_A.
\]
Both conjugates of $\omega,\eta$ are positive. Their smaller conjugates
are $3/\omega<1$ and $1/\eta<1$, so their global heights are
$\tfrac12\log\omega$ and $\tfrac12\log\eta$. Thus
\[
 \height(\gamma^N)\le2\bigl(\height(\omega)+\height(\eta)\bigr)
 <\log(4(b+3))<H+2\le3H.
\]
Applying \eqref{eq:group-approximation} gives
$\log b\ll R_A\log(3H)$, with an effective absolute constant.
Since $H\le\log b+\log3$, we have
$H\le M\log(3H)$ for $M=C_1\max(1,R_A)\ge1$ and an absolute $C_1$.

For any positive $X,M$ with $X\le M\log(3X)$, one has
\begin{equation}\label{eq:log-absorption}
 X\le2M\log(3M).
\end{equation}
To prove it, write $t=X/M$ and use
$\log t\le t/2$, which follows from
$\log(t/2)\le t/2-1$ and $\log2<1$.
Then $t\le\log(3M)+t/2$, proving \eqref{eq:log-absorption}.
Equation~\eqref{eq:quadratic-regulator} bounds $M$ by
$C_2\sqrt A\log(2A)$ and $\log(3M)$ by $C_3\log(2A)$.
Applying \eqref{eq:log-absorption} proves the theorem.
\end{proof}

Replacing $\omega$ by $s+v\sqrt B$, of norm two, gives the analogous
bound with $\sqrt B\log^2(2B)$. For this endpoint, however, a stronger
published input is available and was already identified in the repository.
The unit
\[
 \varepsilon_B=(b+2)+vs\sqrt B
\]
is the fundamental positive norm-one unit: its first coordinate is $3r_0^2$,
so Bennett--Walsh \cite[Theorem~1.2]{BW} says that its unit index is the
first index at which three divides the first Pell coordinate. If the
fundamental first coordinate is $1$ or $-1$ modulo three, the Chebyshev
recurrence gives only $1$ or alternating signs modulo three, respectively.
A first coordinate divisible by three therefore forces the fundamental first
coordinate itself to be divisible by three, and the index is one.
Since $B\equiv3\pmod4$, the order $\Z[\sqrt B]$ is maximal and a
norm-minus-one unit is impossible modulo four. Consequently
\begin{equation}\label{eq:fundamental-B}
 H<R_B\le C_Q\sqrt B\log(2B).
\end{equation}
We use this established endpoint result, rather than presenting it as a
new consequence of the recent preprint.

\begin{corollary}\label{cor:joint-sixth}
There is an effective absolute $K\ge1$ such that these packets satisfy
\begin{equation}\label{eq:joint-sixth}
 H^4\le K D[\log(2D)]^6.
\end{equation}
In particular,
\begin{equation}\label{eq:joint-inverse}
 D\ge\frac{H^4}{K5^6(\log H)^6},\qquad
 \log D\ge4\log H-6\loglog H-\log(K5^6).
\end{equation}
When the latter lower bound is positive,
\[
 p<\frac{4H}{4\log H-6\loglog H-\log(K5^6)}.
\]
\end{corollary}
\begin{proof}
Square Theorem~\ref{thm:fixed-norm-A} and \eqref{eq:fundamental-B} and
multiply. Since $AB=D/3$ and $A,B\le D$, the result is
\eqref{eq:joint-sixth}. If $D\ge H^4$, the first inequality in
\eqref{eq:joint-inverse} is immediate because $H>2$ and $K\ge1$.
Otherwise $\log(2D)\le5\log H$, giving that inequality from
\eqref{eq:joint-sixth}. Taking logarithms proves the second one.
Finally use $p\log(D+1)<4H$.
\end{proof}

The absolute index bound already excludes an unbounded sequence with
$p\ge3$. The kernel estimate \eqref{eq:joint-sixth} remains useful without
a Chebyshev decomposition, including cases where the product Pell unit
is fundamental. Its polynomial dependence on $H$ does not supply the
radical bound required for abc.

\begin{proposition}[A general split-cubic estimate]\label{prop:split-cubic}
For positive squarefree $d$, an integral solution of
$y^2=d(x^3-x)$ with $x\ge2$ satisfies
\[
 \log x\ll d^{2/3}[\log(2d)]^3
\]
with an effective absolute implied constant.
\end{proposition}
\begin{proof}
Write $x-1=a_0u^2$, $x=b_0v^2$, $x+1=c_0w^2$ with positive squarefree
coefficients. Adjacent coefficients are coprime and
$g=\gcd(a_0,c_0)\in\{1,2\}$, so $a_0b_0c_0=g^2d\le4d$.
No two coefficients are equal: equality for adjacent ones would force
positive squares to differ by one; equality at the ends forces the common
coefficient to divide two and gives the same impossibility or a difference
of squares equal to two. The quadratic fields of the three products
$a_0b_0,b_0c_0,a_0c_0$ are therefore distinct.

The three positive norm-one units
\[
 U=2x-1+2uv\sqrt{a_0b_0},\quad
 V=2x+1+2vw\sqrt{b_0c_0},\quad
 W=x+uw\sqrt{a_0c_0}
\]
satisfy $1<V/U,V/(2W),2W/U<1+2/x$.
Indeed $2t-1<t+\sqrt{t^2-1}<2t$ for $t>1$ gives these inequalities
directly at $t=2x-1,2x+1,x$.
Each ratio is generated by two quadratic fundamental units and two;
two may also be adjoined for the first ratio. Distinct quadratic fields
give independence, and two is an independent nonunit. The global height
of each ratio is at most $\log x+O(1)$.
Choose the pair whose radicands share
$m=\min(a_0,b_0,c_0)$. Landau's bound gives a product of regulators at most
\[
 C\sqrt{a_0b_0c_0m}\,[\log(2a_0b_0c_0)]^2
 \le C(a_0b_0c_0)^{2/3}[\log(2a_0b_0c_0)]^2.
\]
Passing to squarefree radicands can only decrease the discriminant bound.
Apply \eqref{eq:group-approximation} to the selected ratio and then
\eqref{eq:log-absorption}, adjusting absolute constants for bounded $x$.
The latter step adds one logarithmic power. Using $a_0b_0c_0\le4d$
proves the claim.
\end{proof}
